
%
%
%
%
%
%
% Alternative Chapter: "Discussion"
\chapter{Discussion: Mechanizing Semantic Proofs — Findings and Insights}

This chapter evaluates the mechanization of semantic proofs in Lean 4 by synthesizing the case-specific insights from the method chapter. The thesis examines how pen-and-paper proofs must be transformed when formalized, with particular attention to implicit assumptions that become explicit, additional machinery required, and lessons learned throughout the mechanization process.

%
%
%
\section{From Informal to Formal: The Transformation of Semantic Definitions}
\label{discussion:sec:transformation}

A primary theme emerging from the formalization of syntactic and semantic
definitions is the systematic conversion of implicit conventions to explicit
requirements.

\subsection*{Syntactic Categories: From BNF to Inductive Types}

The translation of BNF notation to Lean's inductive types is largely
mechanical but reveals several implicit assumptions in textbook presentations:

\begin{itemize}
    \item \textbf{Constructive Numerals}:
    \todo{Fix this, literature does define Num as binary numbers,
    but definition is not needed, could have been Lean 4's Nat Type}
    Textbooks treat numerals as given
    primitives; mechanization requires explicit inductive definition,
    forcing the choice of representation (in our case, binary).

    \item \textbf{Naming and Scoping}:
    Informal notation permits clash-free
    use of keywords like \texttt{if} and \texttt{while}. Mechanization
    requires renaming (\texttt{cond}, \texttt{loop}) to avoid collisions
    with the host language.

    \item \textbf{Boolean Realism}:
    Choosing between symbolic notation and
    computational booleans (using Lean's native \texttt{true}/\texttt{false})
    affects how theorems about boolean logic can be leveraged, reducing
    proof effort.
\end{itemize}

\subsection*{Semantic Functions: Totality and Exhaustiveness}

Formalizing evaluation functions \(\mathcal{N}\), \(\mathcal{A}\), and
\(\mathcal{B}\) reveals that:

\todo{Fix this, literature does define functions as inductive types,
then prove totality, We directly define them as functions to avoid proof,
and to allow application of functions during tests, instead of applying proofs.}

\begin{itemize}
    \item \textbf{Explicit Totality}:
    Lean enforces total function definitions via exhaustive pattern matching.
    Textbooks often implicitly assume totality or leave error cases unstated.
    \todo{Not true, book defines via inductive definition, then proves totality,
    or it defines function as partial as seen with the \(\mathcal{S}\) functions.}
    
    \item \textbf{Constructor Coverage}:
    Each inductive case must be handled explicitly in Lean, whereas textbooks
    omit ``obvious'' or degenerate cases.
    \todo{Give example (such as proving if_false contradicts the b = true case)}
    
    \item \textbf{Type Consistency}:
    Mechanization requires explicit type signatures and consistent return types,
    forcing precision that pen-and-paper notation glosses over.
\end{itemize}

%
%
%
\section{The Explicit Nature of Inductive Proofs}
\label{discussion:sec:inductive_proofs}

Beyond definitions, the mechanization of proofs reveals a systematic pattern:
textbook arguments that appeal to structure
(``the shape of the derivation tree'') must be formalized as explicit case
analysis on inductive constructors.

\subsection*{Case Analysis and the Shape of Derivations}

In natural semantics proofs (such as determinism), the textbook phrase
``the only rule that could have been applied is...'' becomes an explicit
\texttt{cases} tactic in Lean. This forces:

\begin{itemize}
    \item \textbf{Exhaustive Case Analysis}:
    All possible rules must be considered, including those that lead to
    contradictions.
    \todo{Give example (such as proving if_false contradicts the b = true case)}
    \todo{Is this the same bullet point as given in previous section?}
    
    \item \textbf{Side Condition Discharge}:
    Conditions like $\mathcal{B}[\![b]\!]s = \mathbf{tt}$ that the textbook
    uses to rule out certain cases must be explicitly shown to lead to contradictions when negated.
    \todo{Side conditions limit which cases fit the construction,
    lean 4 does not have a counterpart}
    
    \item \textbf{Implicit-to-Explicit Substitution}:
    When the proof uses facts like ``$s_0 = s_1$ therefore...'', mechanization
    requires an explicit \texttt{subst} tactic to propagate the equality
    through all assumptions.
\end{itemize}

\subsection*{Modularity Through Lemmatization}

The while-loop unrolling proof demonstrates that splitting a biconditional
proof into two directed implications (\texttt{while\_expand} and
\texttt{if\_compact}) mirrors good mathematical practice. In mechanization:

\begin{itemize}
    \item \textbf{Independent Lemmas}:
    Proving $A \implies B$ and $B \implies A$ separately, then combining via
    \texttt{Iff.intro}, mirrors the textbook's two-part structure.
    \item \textbf{Constructive Proof Objects}:
    Building derivation trees explicitly (via \texttt{have} statements) makes
    the proof structure transparent in a way that narrative explanations may
    obscure.
    \todo{Using have is not to make the proof more clear than the pen-and-papper,
    but to reorganize the steps taken to create the proof, This is used to move
    a part of the proof earlier in the code to match with the textbook.
    It can however be useful to show a reader our intent}
\end{itemize}

%
%
%
\section{Multi-Step Reasoning and the Abstraction Layers}
\label{discussion:sec:multistep}

A subtle but pervasive difference emerges when formalizing proofs involving
multi-step derivations in structural operational semantics:

\subsection*{From Implicit to Explicit Multi-Step Relations}

The textbook introduces shorthand notation for multiple steps
($\Rightarrow^k$ for $k$ steps, $\Rightarrow^*$ for arbitrary steps).
Mechanization requires:

\begin{itemize}
    \item \textbf{Inductive Types for Relations}:
    The notation $\Rightarrow^k$ must be formalized as an inductive type
    \texttt{small\_step\_k} that counts steps explicitly. This is not merely
    notation but a proof-bearing object.
    
    \item \textbf{Induction on Length}:
    Proofs by ``induction on the length of the derivation sequence'' become
    \texttt{Nat.strongRecOn} in Lean, enabling strong induction rather than
    weak induction.
    
    \item \textbf{Composition of Derivations}:
    The textbook's informal chaining (``we have already argued...'') becomes
    explicit calls to transitivity lemmas like \\
    \texttt{small\_step\_star\_trans}.
\end{itemize}

\subsection*{Case Study: Composition Decomposition (Lemma 2.19)}

The proof that a composition can be split into two independent executions
illustrates how mechanization requires:

\begin{itemize}
    \item \textbf{Automation vs. Manual Proof}:
    Arithmetic facts (e.g., $k_1 + k_2 = k_0 + 1$) are discharged in textbooks
    as ``obvious''. Lean requires \texttt{linarith} or manual simplification.
    
    \item \textbf{Generalization Strategy}:
    To preserve the induction hypothesis across varying statements and states,
    \texttt{generalizing} keywords must be used strategically,
    a concerninvisible in informal mathematics.
    
    \item \textbf{Unpacking and Specialization}:
    Existential witnesses from the induction hypothesis must be unpacked via
    pattern matching and then specialized to the concrete case at hand.
\end{itemize}

%
%
%
\section{Equivalence Proofs: Bridging Two Semantic Styles}
\label{discussion:sec:equivalence}

The two-lemma proof of equivalence between natural semantics and
structural operational semantics ($\texttt{ns\_to\_sos}$ and
$\texttt{sos\_k\_to\_ns}$) demonstrates the deepest level of mechanization
challenges.
\todo{Mention that the book proves equality between partial functions,
but seeing lean 4 has a hard time to reason with partial functions (NOT IMPOSSIBLE)
it is easier to use the direct under the hood meaning instead, as the partial proof
would still need to prove ($\texttt{ns\_to\_sos}$ and $\texttt{sos\_k\_to\_ns}$)
in any case.}

\subsection*{Lemma 2.27: Natural Semantics to Structural Semantics}

This proof constructs multi-step SOS derivations from big-step NS derivations.
Key observations:

\begin{itemize}
    \item \textbf{Constructor-Based Induction}:
    Induction on the shape of the NS derivation tree becomes pattern matching
    on NS inductive constructors (\texttt{ass}, \texttt{skip}, \texttt{comp},
    \texttt{if\_true}, \texttt{if\_false}, \texttt{while\_true},
    \texttt{while\_false}).
    
    \item \textbf{Helper Lemmas}:
    The textbook's informal appeal to ``Exercise 2.21'' becomes an explicit
    function call \texttt{exercise\_2\_21}. Bridging lemmas are no longer
    hidden in prose.
    
    \item \textbf{Transitivity and Composition}:
    Building up multi-step derivations explicitly requires knowing when to
    use \texttt{small\_step\_star.step} versus \texttt{small\_step\_star\_trans},
    a decision the textbook leaves to the reader.
\end{itemize}

\subsection*{Lemma 2.28: Structural Semantics to Natural Semantics}

This proof converts SOS derivations back to NS using strong induction on
step count. Critical differences from Lemma 2.27:

\begin{itemize}
    \item \textbf{Strong vs. Weak Induction}:
    The textbook's informal ``induction on the length'' becomes
    \texttt{Nat.strongRecOn}, enabling the induction hypothesis to apply to
    any smaller step count, not just the predecessor.
    
    \item \textbf{Composition Rule Splitting}:
    The two composition rules [comp$_{\text{sos}}^1$] and
    [comp$_{\text{sos}}^2$] are explicitly split into cases,
    each requiring separate application of Lemma 2.19 to decompose the
    overall derivation.
    
    \item \textbf{Arithmetic Justification}:
    Arguments like ``$k_2 > 0$ because...'' require explicit case analysis
    on natural numbers and use of \texttt{linarith} to discharge arithmetic goals.
    
    \item \textbf{Bridging the While Loop}:
    The while loop rule unrolls into a conditional-composition structure,
    requiring an explicit bridge via \texttt{lemma\_2\_5} to convert back
    to the while rule in NS.
\end{itemize}

\subsection*{Completing the Equivalence}

The final theorem \texttt{ns\_equvivlent\_sos} shows that biconditional
proofs require:

\begin{itemize}
    \item \textbf{Relational vs. Functional Interpretation}:
    While the textbook states ``the semantic functions are equal'',
    mechanization works with the underlying evaluation relations.
    This is not a choice but a necessity: partial functions cannot
    directly be used in Lean proofs.
    \todo{It is not directly necessary, as we can model partial functions,
    in this case there is no gain, and we simply add one more step to prove
    what we already have.}
    
    \item \textbf{Biconditional Construction}:
    Using \texttt{Iff.intro} to explicitly bind the forward and backward
    implications mirrors the textbook's two-lemma structure, but makes
    the dependency explicit.
\end{itemize}

%
%
%
\section{Implicit Becomes Explicit: Key Mechanization Challenges}
\label{discussion:sec:implicit_explicit}

Across all proof cases, a consistent pattern emerges: pen-and-paper
proofs rely on informal conventions and reader understanding that
mechanization must formalize.

\subsection*{Totality and Partiality}

Textbooks often elide the distinction between total and partial functions.
Mechanization requires:

\begin{itemize}
    \item \textbf{Total Functions}:
    Evaluation functions like \(\mathcal{A}\) must be total,
    enforced by exhaustive pattern matching.
    
    \item \textbf{Partial Semantic Functions}:
    The semantic function \(\mathcal{S}\)
    (from statement and initial state to final state)
    is inherently partial: not all programs terminate.
    This cannot be directly represented in Lean's type system; instead,
    the evaluation relation is used as a more fundamental primitive.
    \todo{Not true, see `Thesis/Definitions/Partial.lean`
    for a partial implementation of the functions.}
\end{itemize}

\subsection*{Side Conditions and Disjunctive Cases}

Textbook rules often include side conditions
(e.g., $\mathcal{B}[\![b]\!]s = \mathbf{tt}$).
Mechanization requires:

\begin{itemize}
    \item \textbf{Explicit Proof of Conditions}:
    Every application of a rule must discharge the side condition as a goal.
    
    \item \textbf{Contradiction Checking}:
    When a case is ruled out (e.g., when the condition is false and the rule
    requires it to be true), mechanization explicitly derives a
    contradiction (e.g., via \texttt{rw} and \texttt{contradiction}).
\end{itemize}

\subsection*{Witness and Generalization}

Existential quantification and generalization introduce subtle mechanization
burdens:

\begin{itemize}
    \item \textbf{Existential Witnesses}:
    Unpacking witnesses from the induction hypothesis requires explicit
    pattern matching, followed by specialization to the concrete instance.
    
    \item \textbf{Generalization in Induction}:
    Variables that should not be fixed during induction
    (e.g., the second derivation $s''$ in the determinism proof)
    must be explicitly generalized via the \texttt{generalizing} keyword.
\end{itemize}

\subsection*{Arithmetic and Automation}

Pen-and-paper proofs often dismiss arithmetic as ``obvious'' or ``trivial''.
Mechanization requires:

\begin{itemize}
    \item \textbf{Tactic Automation}:
    Lemmas like \texttt{simp} and \texttt{linarith} handle arithmetic
    bookkeeping, but their application must be explicit.
    
    \item \textbf{Manual Justification}:
    When automation suffices, the proof is concise. When it doesn't,
    manual cases (e.g., \texttt{cases k₂}) and lemmas become necessary.
\end{itemize}

%
%
%
\section{Proof Techniques That Bridged Theory and Practice}
\label{discussion:sec:techniques}

Several proof strategies proved particularly effective in bridging
the gap between textbook and mechanized proofs.

\subsection*{Induction on the Shape of Derivation Trees}

Natural semantics proofs use ``induction on the shape of the derivation tree''.
In mechanization:

\begin{itemize}
    \item \textbf{Tactic: \texttt{induction}}:
    Lean's \texttt{induction} tactic directly mirrors this pattern,
    automatically generating cases for each inductive constructor.
    
    \item \textbf{Case Naming}:
    Automatically generated case names correspond to rule labels
    (e.g., \texttt{case ass}, \texttt{case comp}), making the correspondence
    with textbook rules transparent.
\end{itemize}

\subsection*{Strong Induction on Derivation Length}

SOS proofs use ``induction on the length of the derivation sequence''.
In mechanization:

\begin{itemize}
    \item \textbf{Tactic: \texttt{Nat.strongRecOn}}:
    Strong induction enables the induction hypothesis to apply to any
    smaller step count, essential for splitting a multi-step derivation
    into sub-derivations of unknown length.
    
    \item \textbf{Proof by Cases}:
    Pattern matching on the natural number (zero vs. successor)
    explicitly captures the base and inductive cases.
\end{itemize}

\subsection*{Lemmatization and Reuse}

Breaking complex proofs into lemmas (e.g., \texttt{seq\_split},
\texttt{exercise\_2\_21}) enables:

\begin{itemize}
    \item \textbf{Modular Reasoning}:
    Each lemma can be understood and verified independently.
    
    \item \textbf{Proof Reuse}:
    Lemmas like \texttt{seq\_split} are applied multiple times,
    formalizing the textbook's informal appeal to the same reasoning in
    different contexts.
\end{itemize}


%
%
%
\section{Lessons for Future Mechanizations}
\label{discussion:sec:lessons}

Based on the mechanization experience,
several lessons emerge for future work on semantic formalization.

\subsection*{Planning and Structuring Proofs}

\begin{itemize}
    \item \textbf{Identify Proof Patterns Early}:
    Recognize whether the proof is structural (shape of derivations)
    or arithmetic (length of sequences).
    Choose induction strategies accordingly.
    
    \item \textbf{Lemmatize Aggressively}:
    Large proofs should be broken into smaller,
    reusable lemmas. This mirrors textbook structure but also improves
    proof maintainability.
    
    \item \textbf{Use Helper Lemmas Explicitly}:
    What appears as informal appeal in a textbook should
    become explicit lemma calls in mechanization.
\end{itemize}

\subsection*{Handling Implicit Assumptions}

\begin{itemize}
    \item \textbf{Side Conditions}:
    Enumerate all side conditions and plan how to discharge them.
    Cases that are ``ruled out'' by side conditions must be shown to
    lead to contradictions.
    
    \item \textbf{Variable Generalization}:
    When setting up induction, generalize any variable that should
    not be fixed by the induction. This is often a source of subtle
    proof failures.
    
    \item \textbf{Existential Witnesses}:
    Plan how to unpack and specialize existential witnesses.
    Pattern matching is powerful but requires careful ordering.
\end{itemize}

\subsection*{Arithmetic and Automation}

\begin{itemize}
    \item \textbf{Rely on Automation Where Possible}:
    Lemmas like \texttt{simp} and \texttt{linarith} handle routine bookkeeping;
    use them liberally.
    
    \item \textbf{Separate Concerns}:
    Keep non-arithmetic reasoning separate from arithmetic reasoning.
    This makes both easier to verify independently.
    
    \item \textbf{Document Automation Justification}:
    When relying on \texttt{linarith} or \texttt{simp},
    a comment explaining the intended simplification helps future readers.
\end{itemize}

\subsection*{The Semantic Function Dilemma}

\begin{itemize}
    \item \textbf{Work with Relations, Not Functions}:
    Partial functions (like semantic evaluation from statement and state
    to final state) cannot be directly formalized in dependent type theory.
    Instead, work with the evaluation relation as the primary object
    and derive the semantic function as a derived notion if needed.

    \item \textbf{Equivalence of Relations is Practical}:
    Equivalence of evaluation relations is more tractable and sufficient
    for most applications.
\end{itemize}

\subsection*{Collaboration with Textbooks}
\todo{What is the purpose of this part?}
\begin{itemize}
    \item \textbf{Align Notation Early}:
    Choose mechanized notation that mirrors the textbook closely.
    This eases comparison and reader comprehension.
    
    \item \textbf{Side-by-Side Presentation}:
    Present mechanized and textbook proofs side by side during
    development to catch divergences early.
    
    \item \textbf{Document Deviations}:
    When mechanization differs from the textbook
    (e.g., using strong induction instead of weak),
    document the reason and implications.
\end{itemize}

%
%
%
\section{Synthesis: Toward a Typology of Mechanization Differences}
\label{discussion:sec:synthesis}

Synthesizing all insights,
the systematic differences between pen-and-paper and
mechanized proofs fall into several categories:

\begin{table}[h]
    \centering
    \begin{tabular}{p{3cm}|p{4cm}|p{4cm}}
        \hline
        \textbf{Category} & \textbf{Pen-and-Paper} & \textbf{Mechanized} \\
        \hline
        \textit{Totality} & Often implicit & Enforced by exhaustive matching \\
        \hline
        \textit{Side Conditions} & Ruled out informally & Treated as a premise that can lead to contradictions \\
        \hline
        \textit{Case Analysis} & Narrative description & Explicit \texttt{cases} tactic \\
        \hline
        \textit{Multi-Step Reasoning} & Implicit chaining & Explicit transitivity lemmas \\
        \hline
        \textit{Generalization} & Implicit scope & Explicit \texttt{generalizing} keywords \\
        \hline
        \textit{Arithmetic} & Dismissed as obvious & Handled by automation or manual cases \\
        \hline
        \textit{Lemma Application} & Prose appeal & Explicit function calls with arguments \\
        \hline
        \textit{Witness Unpacking} & Informal existential & Pattern matching with specialization \\
        \hline
    \end{tabular}
\end{table}

This typology suggests that mechanization is not merely a formal restatement of textbook arguments, but a systematic transformation that makes every implicit assumption explicit and every logical step verifiable.

\todo{Maybe it is worth mentioning that proofs written in Lean 4 can have many shapes
and what is provided is made to mimic the textbook as close as possible,
This is important as compared to written proofs, knowing a what a proof is proving,
and that it is compiled in enough to verify the correctness, whilst text based proofs
need manual verification of every individual step to verify no errors have occured.}